% !TEX program = xelatex
\documentclass[aspectratio=169]{ctexbeamer}
% 编译方式：xelatex Fisher-R1-Report.tex

\usetheme{Madrid}
\usecolortheme{whale}
\setbeamertemplate{navigation symbols}{}
\setbeamerfont{footline}{size=\small}

\usepackage{booktabs}
\usepackage{amsmath, amssymb, amsfonts}
\usepackage{array}

\title[Fisher-R1 汇报]{Fisher-R1: Training LLM Agents\\for Reliable Hypothesis Testing}
\subtitle{可执行的分析 $\neq$ 可靠的统计推断}
\author{汇报人}
\date{2026 年 8 月}

\begin{document}

% ============================================================
\begin{frame}
  \titlepage
\end{frame}

% ============================================================
\begin{frame}{Outline}
  \tableofcontents
\end{frame}

% ============================================================
\section{Motivation: 现有 LLM Agent 的失败模式}

\begin{frame}{Opening: 一张图看懂全文（Figure 1 的故事）}
  \textbf{场景}：数据中存在两个 \alert{high-leverage outliers}，同样的数据：
  \begin{itemize}
    \item Linear regression 给出 $p < 2.2\times 10^{-16}$（极显著）
    \item Spearman rank test 给出 $p = 0.085$（不显著）
  \end{itemize}

  \medskip
  \textbf{GPT-5.4 的表现}：
  \begin{itemize}
    \item 看到了 outlier
    \item 算过 Spearman
    \item 知道 linear regression 有 warning signs
    \item \alert{最终仍然选择 linear regression}
  \end{itemize}

  \medskip
  \textbf{Fisher-R1}：最终选择 rank test。

  \medskip
  \begin{beamercolorbox}[center, sep=8pt]{box}
    \Large Correct execution $\neq$ valid statistical inference
  \end{beamercolorbox}
\end{frame}

\begin{frame}{为什么需要专门的 Benchmark？}
  \begin{itemize}
    \item $p$-value 依赖于所选的 statistical model；错误的 method selection
          导致错误的 conclusion。
    \item 现有通用模型（GPT-5.4 等）可以可靠地看数据、生成代码、执行分析、
          输出报告——\textbf{但所选统计模型本身未必可靠}。
    \item 现有 benchmark 很少评估 ``报告的 $p$-value 在数据假设下是否统计有效''，
          无法捕捉这一 failure mode。
  \end{itemize}

  \medskip
  \begin{beamercolorbox}[center, sep=8pt]{box}
    汇报重点排序：Motivation $>$ P-Bench $\approx$ Synthetic Verifiable Training
    $>$ Experimental Findings $>$ Fisher Reward $>$ DAPO Details
  \end{beamercolorbox}
\end{frame}

% ============================================================
\section{P-Bench（一）：它到底测什么}

\begin{frame}{P-Bench 评测的完整链条}
  \textbf{输入}：一个科学问题的分析请求 + 一个 CSV 数据集。
  \alert{不告知模型应使用何种统计方法。}

  \medskip
  模型需要独立完成完整链条：
  \[
    \boxed{\;\text{question} \rightarrow \text{data exploration}
      \rightarrow \text{method selection} \rightarrow \text{execute R}
      \rightarrow p\text{-value} \rightarrow \text{decision}\;}
  \]

  \medskip
  \textbf{与 StatQA 类选择题 benchmark 的核心区别}：
  \begin{itemize}
    \item 不是从选项中挑方法，而是 end-to-end 自主决策 + 执行；
    \item 评的不只是 ``方法选对没''，还有 $p$-value 与 conclusion 的统计有效性。
  \end{itemize}
\end{frame}

% ============================================================
\section{P-Bench（二）：构造与 Ground Truth 的可信性}

\begin{frame}{425 个任务从哪里来？}
  \[
    203\ \text{Easy} + 222\ \text{Hard} = \boxed{425}
  \]

  \begin{itemize}
    \item 来自经济学、生物学、医学领域的\textbf{真实论文}及其数据集；
    \item 覆盖 \textbf{17 类}实际使用的 hypothesis-testing methods：
      \begin{itemize}
        \item 常规方法：$t$-test、$\chi^2$、OLS、Fisher exact 等
        \item 高级方法：Cox PH、IV/2SLS、mixed effects、Tobit 等
      \end{itemize}
    \item Hard 不仅是 ``题目更复杂''，而是包含 statistical traps：
      model specification、assumption checking、outliers、
      heteroskedasticity、clustered observations。
  \end{itemize}
\end{frame}

\begin{frame}{Ground truth 为什么可信？（关键一页）}
  答案不是从论文文本里抽一个 $p$-value，而是来自
  \alert{canonical analysis 的真实执行}：
  \[
    \text{paper + dataset + analysis code}
    \;\rightarrow\; \text{重新运行}
    \;\rightarrow\; \text{execution log}
    \;\rightarrow\; \text{提取 } p^{*},\ \text{conclusion}
    \;\rightarrow\; \text{专家审核}
  \]

  \medskip
  \begin{itemize}
    \item Keys 由 canonical reference code 计算得到，并经专家检查。
    \item 应对质疑：``你说 GPT-5.4 错，凭什么 benchmark 的答案是对的？''
          ——因为答案是\textbf{可执行、可复现的 canonical 代码}算出来的，
          而非模型或人工的主观判断。
  \end{itemize}
\end{frame}

% ============================================================
\section{Fisher-R1（一）：六维 Synthetic Task Space}

\begin{frame}{Task Taxonomy：合成任务不是乱生成的}
  \[
    \boxed{\;M \times S_m \times N \times E \times P \times K\;}
  \]
  \begin{center}
  \begin{tabular}{c l}
    \toprule
    维度 & 含义 \\
    \midrule
    $M$   & statistical method \\
    $S_m$ & scenario（应用场景） \\
    $N$   & sample size \\
    $E$   & effect size（null / borderline / medium） \\
    $P$   & prompt style（hint / non-hint） \\
    $K$   & random seed \\
    \bottomrule
  \end{tabular}
  \end{center}

  \medskip
  附录训练空间：\textbf{27 种} statistical methods、每种 5--7 个场景、
  多种 $N$ 与 effect size、hint/non-hint prompt，共合成 \textbf{8642} 个任务。
\end{frame}

\begin{frame}{为什么 Synthetic 让 RL 成为可能}
  每一个坐标点产生一个\textbf{可执行}的统计任务，自然得到
  machine-verifiable answer key：
  \[
    \text{DGP} \rightarrow \text{dataset.csv}
    \qquad\qquad
    \text{dataset} \xrightarrow{\ \text{canonical method}\ } p^{*}
  \]

  \medskip
  \begin{beamercolorbox}[center, sep=8pt]{box}
    核心创新不是 ``我们用了 DAPO 训练 Qwen''，而是\\
    \textbf{把一个难以自动监督的统计推断问题，转化为可自动生成
    ground truth、可自动计算 reward、进而可做 RL 的任务。}
  \end{beamercolorbox}

  \medskip
  方法主线压缩为：
  \[
    \text{Synthetic executable tasks}
    \rightarrow \text{expert SFT warm start}
    \rightarrow \text{outcome-verifiable RL}
  \]
\end{frame}

% ============================================================
\section{Fisher-R1（二）：SFT Warm Start}

\begin{frame}{SFT：轻量带过}
  \begin{itemize}
    \item Teacher：Claude Sonnet 4.6；
    \item 固定\textbf{五步工作流}生成轨迹：
      \[
        \text{EDA} \rightarrow \text{detailed EDA}
        \rightarrow \text{assumption checking}
        \rightarrow \text{method selection}
        \rightarrow \text{conclusion}
      \]
    \item 共生成 4611 条 trajectory，筛选后保留 3851 条；
    \item 格式化为 ReAct / CodeAct 格式；
    \item 作为 Qwen（7B / 14B）的 SFT warm-start。
  \end{itemize}

  \medskip
  \textit{备注（若被问 SFT loss）}：prompt / observation 只作为 context，
  只有 assistant tokens 计算 loss。
\end{frame}

% ============================================================
\section{Fisher-R1（三）：Outcome-Verifiable RL}

\begin{frame}{Fisher Reward（重点）}
  先回答：\textbf{为什么这里能做 RL？}——因为 synthetic task 自带
  verifiable ground truth $p^{*}$。于是定义（论文原式）：
  \[
    R(o) = I_{\mathrm{valid}}(o)\left[w_p\, r_p(o) + w_c\, r_c(o)\right],
    \qquad w_p = 0.9,\quad w_c = 0.1
  \]

  \medskip
  \textbf{90\%：$p$-value accuracy}。先做 $z$ 变换再比较：
  \[
    z(p) = \Phi^{-1}\!\left(1 - \tfrac{p}{2}\right),
    \qquad
    r_p(o) = \exp\!\left(
      -\frac{\bigl|\min\{z(\hat p),5\} - \min\{z(p^{*}),5\}\bigr|}{\sigma}
    \right)
  \]
  为什么不用 raw $p$-difference？$0.5\leftrightarrow 0.6$ 与
  $0.1\leftrightarrow 10^{-10}$ 的 $|\Delta p|$ 可能接近，
  但 statistical evidence 的差距完全不同。

  \medskip
  \textbf{10\%：conclusion accuracy}。判断 reject / fail-to-reject 是否正确。
\end{frame}

\begin{frame}{DAPO：两分钟讲完}
  They optimize this reward using \textbf{DAPO},
  a group-relative policy optimization method：
  \[
    \hat A_i = \frac{R_i - \mu}{\sigma + \epsilon}
  \]
  \begin{itemize}
    \item $\hat A_i$ 代表\textbf{同一道题}多个 rollout 中的相对质量；
    \item positive advantage $\uparrow$ probability，
          negative advantage $\downarrow$ probability；
    \item clip 保证单步更新不要过于激进；
    \item 丢弃 $\sigma_q^2 = 0$ 的 group——即同组所有 rollout 的 reward
          \alert{没有差异}（可能全 0、全 1 或全部相同值），
          目的是丢掉\textbf{没有 relative advantage learning signal} 的组，
          而非专门丢弃失败组。
  \end{itemize}
\end{frame}

% ============================================================
\section{实验结果}

\begin{frame}{Main Results（Table 1）：只挑三个事实}
  \textbf{事实 1：post-training 提升巨大}
  \[
    7\text{B backbone: } 13.2\% \rightarrow 30.6\%
    \quad (\text{Hard Strict pass@1})
  \]

  \medskip
  \textbf{事实 2：14B 部分超越 GPT-5.4}\\
  Fisher-R1-14B 在四个 Strict metric 中三个超过 GPT-5.4，例如 Hard：
  \[
    33.0\% \ \text{ vs. }\ 30.5\%
  \]
  \textit{（注意分寸：不是做成 ``我们打赢了 GPT-5.4'' 的叙事。）}

  \medskip
  \textbf{事实 3：Raw vs.\ Strict gap（最具科学意义）}
  \[
    \text{GPT-5.4 Hard pass@1: } \text{Raw} = 58.3\%,\quad
    \text{Strict} = 30.5\%
  \]
  直接支持论文 thesis：模型常能得到正确的显著方向，
  但无法可靠地产生正确的 statistical evidence。
\end{frame}

\begin{frame}{Raw vs.\ Strict Gap 的含义}
  \begin{beamercolorbox}[center, sep=10pt]{box}
    \Large Conclusion accuracy can substantially overestimate\\
    statistical reliability.
  \end{beamercolorbox}

  \medskip
  \begin{itemize}
    \item Raw（结论方向对）$\neq$ Strict（$p$-value 与方法也对）；
    \item 这正是 Figure 1 失败模式的定量版本：执行正确、结论方向碰对，
          但 statistical evidence 无效；
    \item 也因此 P-Bench 的 Strict 评价才是这篇工作真正的贡献点。
  \end{itemize}
\end{frame}

\begin{frame}{Ablation \& Generalization（一页）}
  \textbf{Ablation（P-Hard Strict pass@1, \%）}：
  \begin{center}
  \begin{tabular}{l c}
    \toprule
    设置 & pass@1 \\
    \midrule
    Backbone only          & 13.2 \\
    + DAPO only            & 25.2 \\
    + SFT only             & 14.3 \\
    SFT + DAPO             & \textbf{30.6} \\
    \bottomrule
  \end{tabular}
  \end{center}
  \[
    \boxed{\text{SFT warm-start 与 outcome RL 是互补的。}}
  \]

  \medskip
  \textbf{Generalization}：训练任务全是 synthetic、P-Bench 任务是真实的；
  作者做了 embedding nearest-neighbor 分析，未发现 P-Bench 任务与训练
  prompt 存在 near-duplicate。
\end{frame}

% ============================================================
\section{Limitations 与个人讨论}

\begin{frame}{Limitations}
  \begin{itemize}
    \item 目前只评估\textbf{单个} hypothesis test；
    \item 未覆盖 multiple-testing / 多阶段分析 pipeline；
    \item method choice 存在 uncertainty：同一问题常有多个 defensible
          statistical procedures；
    \item 有的问题可能根本不存在单一的 adequate test。
  \end{itemize}
\end{frame}

\begin{frame}{我的 Critique：Method 非唯一 vs.\ Reward 指向唯一}
  作者一方面承认\textbf{多个 statistical procedures 都可能 defensible}，
  因此 reward 不直接奖励 method correctness；另一方面又通过逼近某个
  canonical $p^{*}$ 来评价模型。于是存在 tension：
  \[
    \boxed{\text{method may be non-unique}}
    \qquad\text{but}\qquad
    \boxed{\text{reward target is essentially one canonical inference}}
  \]

  \medskip
  因此 Fisher-R1 更准确的定性是学习：
  \[
    \text{produce inference close to an expert canonical analysis}
  \]
  而不能等价为：
  \[
    \text{discover the uniquely correct statistical method.}
  \]
\end{frame}

% ============================================================
\section{Takeaway}

\begin{frame}{Takeaway}
  \begin{enumerate}
    \setlength{\itemsep}{12pt}
    \item \textbf{Executable analysis does not imply valid inference.}
    \item \textbf{P-Bench} 用可执行、专家级 ground truth 的真实任务，
          暴露并量化了这一 gap。
    \item \textbf{Synthetic verifiable post-training}（SFT + outcome RL）
          可以大幅提升统计推断的可靠性。
    \end{enumerate}

  \medskip
  \textit{Reference: Miao, Mu, Chen \& Zou.
  Fisher-R1: Training LLM Agents for Reliable Hypothesis Testing.
  arXiv:2608.07437, 2026.}
\end{frame}

\begin{frame}
  \begin{center}
    \Huge Thanks!\\[12pt]
    \Large Q \& A
  \end{center}
\end{frame}

\end{document}
